
\documentclass{article}
\usepackage{colortbl}
\usepackage{makecell}
\usepackage{multirow}
\usepackage{supertabular}

\begin{document}

\newcounter{utterance}

\centering \large Interaction Transcript for game `dond', experiment `coop', episode 4 with nemotron{-}3{-}ultra.
\vspace{24pt}

{ \footnotesize  \setcounter{utterance}{1}
\setlength{\tabcolsep}{0pt}
\begin{supertabular}{c@{$\;$}|p{.15\linewidth}@{}p{.15\linewidth}p{.15\linewidth}p{.15\linewidth}p{.15\linewidth}p{.15\linewidth}}
    \# & $\;$A & \multicolumn{4}{c}{Game Master} & $\;\:$B\\
    \hline

    \theutterance \stepcounter{utterance}  
    & & \multicolumn{4}{p{0.6\linewidth}}{
        \cellcolor[rgb]{0.9,0.9,0.9}{
            \makecell[{{p{\linewidth}}}]{
                \texttt{\tiny{[P1$\langle$GM]}}
                \texttt{Grasz w grę negocjacji kooperacyjnej, w której musisz uzgodnić z innym graczem sposób podziału zestawu przedmiotów między sobą.} \\
\\ 
\texttt{Zasady:} \\
\texttt{(a) Ty i drugi gracz otrzymujecie zestaw przedmiotów. Każdemu z Was przydzielana jest również tajna funkcja wartości, reprezentująca, ile warte są dla Ciebie poszczególne typy przedmiotów.} \\
\texttt{(b) Wymieniasz wiadomości z drugim graczem, aby uzgodnić, kto otrzyma jakie przedmioty. Każdy może wysłać maksymalnie 5 wiadomości lub skończyć grę wcześniej, składając tajną propozycję w dowolnym momencie.} \\
\texttt{(c) Poproszony o złożenie tajnej propozycji wskazującej przedmioty, które chcesz, sformatowanej w nawiasach kwadratowych w następujący sposób: "{[}Propozycja: <liczba> <nazwa przedmiotu>, <liczba> <nazwa przedmiotu>, <...>{]}"} \\
\texttt{(d) Jeśli Twoje propozycje są się się nawzajem uzupełniają, tzn. jest wystarczająco przedmiotów, aby spełnić obie propozycje, każdy gracz otrzymuje wynik na podstawie sumy wartości przedmiotów, które otrzymał. W przeciwnym razie obaj otrzymujecie zero punktów.} \\
\texttt{(e) Twój cel to maksymalizacja sumy Twojego wyniku i wyniku drugiego gracza.} \\
\\ 
\texttt{Zacznijmy.} \\
\\ 
\texttt{Zestaw dostępnych przedmiotów to:} \\
\\ 
\texttt{3 książki, 1 torba, 1 śruba} \\
\\ 
\texttt{Twoja tajna funkcja wartości to:} \\
\\ 
\texttt{książka: 0, torba: 2, śruba: 8} \\
\\ 
\texttt{WAŻNE: Twoje wiadomości, chyba że jest to tajna propozycja, są przesyłane bezpośrednio drugiemu graczowi, więc nie dodawaj żadnej odpowiedzi na zasady ani tekstu ogłaszającego Twoją wiadomość. Aby złożyć tajną propozycję, użyj wskazanego formatu. Nie używaj nawiasów kwadratowych podczas komunikowania się z drugim graczem, w przeciwnym razie zostanie to zinterpretowane jako Twoja tajna propozycja.} \\
            }
        }
    }
    & \\ \\

    \theutterance \stepcounter{utterance}  
    & \multicolumn{4}{p{0.6\linewidth}}{
        \cellcolor[rgb]{0.9,0.9,0.9}{
            \makecell[{{p{\linewidth}}}]{
                \texttt{\tiny{[P1$\rangle$GM]}}
                \texttt{Cześć! Widzisz te przedmioty? Dla mnie kluczowa jest śruba (warta 8 punktów), a torba ma wartość 2. Książki nie mają dla mnie żadnej wartości (0 punktów).} \\
\\ 
\texttt{Co Ty zyskujesz na poszczególnych przedmiotach? Jeśli nie chcesz śruby ani torby, mogę je wziąć, a Ty możesz zabrać wszystkie książki – wtedy obaj wygrywamy. Jak Tobie to brzmi?} \\
            }
        }
    }
    & & \\ \\

    \theutterance \stepcounter{utterance}  
    & & & \multicolumn{2}{p{0.3\linewidth}}{
        \cellcolor[rgb]{0.9,0.9,0.9}{
            \makecell[{{p{\linewidth}}}]{
                \texttt{\tiny{[GM$|$GM]}}
                \texttt{continue} \\
            }
        }
    }
    & & \\ \\

    \theutterance \stepcounter{utterance}  
    & & & \multicolumn{2}{p{0.3\linewidth}}{
        \cellcolor[rgb]{0.9,0.9,0.9}{
            \makecell[{{p{\linewidth}}}]{
                \texttt{\tiny{[GM$|$GM]}}
                \texttt{Cześć! Widzisz te przedmioty? Dla mnie kluczowa jest śruba (warta 8 punktów), a torba ma wartość 2. Książki nie mają dla mnie żadnej wartości (0 punktów).} \\
\\ 
\texttt{Co Ty zyskujesz na poszczególnych przedmiotach? Jeśli nie chcesz śruby ani torby, mogę je wziąć, a Ty możesz zabrać wszystkie książki – wtedy obaj wygrywamy. Jak Tobie to brzmi?} \\
            }
        }
    }
    & & \\ \\

    \theutterance \stepcounter{utterance}  
    & & \multicolumn{4}{p{0.6\linewidth}}{
        \cellcolor[rgb]{0.9,0.9,0.9}{
            \makecell[{{p{\linewidth}}}]{
                \texttt{\tiny{[GM$\rangle$P2]}}
                \texttt{Grasz w grę negocjacji kooperacyjnej, w której musisz uzgodnić z innym graczem sposób podziału zestawu przedmiotów między sobą.} \\
\\ 
\texttt{Zasady:} \\
\texttt{(a) Ty i drugi gracz otrzymujecie zestaw przedmiotów. Każdemu z Was przydzielana jest również tajna funkcja wartości, reprezentująca, ile warte są dla Ciebie poszczególne typy przedmiotów.} \\
\texttt{(b) Wymieniasz wiadomości z drugim graczem, aby uzgodnić, kto otrzyma jakie przedmioty. Każdy może wysłać maksymalnie 5 wiadomości lub skończyć grę wcześniej, składając tajną propozycję w dowolnym momencie.} \\
\texttt{(c) Poproszony o złożenie tajnej propozycji wskazującej przedmioty, które chcesz, sformatowanej w nawiasach kwadratowych w następujący sposób: "{[}Propozycja: <liczba> <nazwa przedmiotu>, <liczba> <nazwa przedmiotu>, <...>{]}"} \\
\texttt{(d) Jeśli Twoje propozycje są się się nawzajem uzupełniają, tzn. jest wystarczająco przedmiotów, aby spełnić obie propozycje, każdy gracz otrzymuje wynik na podstawie sumy wartości przedmiotów, które otrzymał. W przeciwnym razie obaj otrzymujecie zero punktów.} \\
\texttt{(e) Twój cel to maksymalizacja sumy Twojego wyniku i wyniku drugiego gracza.} \\
\\ 
\texttt{Zacznijmy.} \\
\\ 
\texttt{Zestaw dostępnych przedmiotów to:} \\
\\ 
\texttt{3 książki, 1 torba, 1 śruba} \\
\\ 
\texttt{Twoja tajna funkcja wartości to:} \\
\\ 
\texttt{książka: 1, torba: 2, śruba: 5} \\
\\ 
\texttt{WAŻNE: Twoje wiadomości, chyba że jest to tajna propozycja, są przesyłane bezpośrednio drugiemu graczowi, więc nie dodawaj żadnej odpowiedzi na zasady ani tekstu ogłaszającego Twoją wiadomość. Aby złożyć tajną propozycję, użyj wskazanego formatu. Nie używaj nawiasów kwadratowych podczas komunikowania się z drugim graczem, w przeciwnym razie zostanie to zinterpretowane jako Twoja tajna propozycja.} \\
\\ 
\texttt{Cześć! Widzisz te przedmioty? Dla mnie kluczowa jest śruba (warta 8 punktów), a torba ma wartość 2. Książki nie mają dla mnie żadnej wartości (0 punktów).} \\
\\ 
\texttt{Co Ty zyskujesz na poszczególnych przedmiotach? Jeśli nie chcesz śruby ani torby, mogę je wziąć, a Ty możesz zabrać wszystkie książki – wtedy obaj wygrywamy. Jak Tobie to brzmi?} \\
            }
        }
    }
    & \\ \\

    \theutterance \stepcounter{utterance}  
    & & & \multicolumn{4}{p{0.6\linewidth}}{
        \cellcolor[rgb]{0.9,0.9,0.9}{
            \makecell[{{p{\linewidth}}}]{
                \texttt{\tiny{[GM$\langle$P2]}}
                \texttt{Dla mnie wartości wyglądają tak: książka 1, torba 2, śruba 5.} \\
\\ 
\texttt{Twój plan brzmi dobrze – śruba i torba dla Ciebie (razem 10 pkt), a dla mnie 3 książki (3 pkt). Suma 13 punktów.} \\
\\ 
\texttt{Zatwierdzam ten podział. Składam propozycję.} \\
\texttt{{[}Propozycja: 3 książki{]}} \\
            }
        }
    }
     \\ \\

    \theutterance \stepcounter{utterance}  
    & & & \multicolumn{2}{p{0.3\linewidth}}{
        \cellcolor[rgb]{0.9,0.9,0.9}{
            \makecell[{{p{\linewidth}}}]{
                \texttt{\tiny{[GM$|$GM]}}
                \texttt{proposal} \\
            }
        }
    }
    & & \\ \\

    \theutterance \stepcounter{utterance}  
    & & & \multicolumn{2}{p{0.3\linewidth}}{
        \cellcolor[rgb]{0.9,0.9,0.9}{
            \makecell[{{p{\linewidth}}}]{
                \texttt{\tiny{[GM$|$GM]}}
                [3, 0, 0]
            }
        }
    }
    & & \\ \\

    \theutterance \stepcounter{utterance}  
    & & \multicolumn{4}{p{0.6\linewidth}}{
        \cellcolor[rgb]{0.9,0.9,0.9}{
            \makecell[{{p{\linewidth}}}]{
                \texttt{\tiny{[P1$\langle$GM]}}
                \texttt{Drugi gracz złożył tajną propozycję. Teraz Twoja kolej na złożenie własnej propozycji.} \\
            }
        }
    }
    & \\ \\

    \theutterance \stepcounter{utterance}  
    & \multicolumn{4}{p{0.6\linewidth}}{
        \cellcolor[rgb]{0.9,0.9,0.9}{
            \makecell[{{p{\linewidth}}}]{
                \texttt{\tiny{[P1$\rangle$GM]}}
                \texttt{{[}Propozycja: 1 torba, 1 śruba{]}} \\
            }
        }
    }
    & & \\ \\

    \theutterance \stepcounter{utterance}  
    & & & \multicolumn{2}{p{0.3\linewidth}}{
        \cellcolor[rgb]{0.9,0.9,0.9}{
            \makecell[{{p{\linewidth}}}]{
                \texttt{\tiny{[GM$|$GM]}}
                \texttt{proposal} \\
            }
        }
    }
    & & \\ \\

    \theutterance \stepcounter{utterance}  
    & & & \multicolumn{2}{p{0.3\linewidth}}{
        \cellcolor[rgb]{0.9,0.9,0.9}{
            \makecell[{{p{\linewidth}}}]{
                \texttt{\tiny{[GM$|$GM]}}
                [0, 1, 1]
            }
        }
    }
    & & \\ \\

    \theutterance \stepcounter{utterance}  
    & & & \multicolumn{2}{p{0.3\linewidth}}{
        \cellcolor[rgb]{0.9,0.9,0.9}{
            \makecell[{{p{\linewidth}}}]{
                \texttt{\tiny{[GM$|$GM]}}
                [[0, 1, 1], [3, 0, 0]]
            }
        }
    }
    & & \\ \\

\end{supertabular}
}

\end{document}
